\section{Conclusions}\label{sec:boundaries}
We classify exact shared continuous planar encoders for the desired functions $f_{\phi,\psi}(\mathbf s)=t-\psi(m)$ and optimize their worst-case conditioning under pointwise peak power. We turn this structural result into a sharp uniform disturbance guarantee by combining the exact score, convex encoder design, and projection decoder. The inverse-metric description also yields the global robustness-constrained small-noise frontier for the specified decoder. Its quadratic specialization quantifies the cost of adapting to an input distribution, while certified discretization supplies a stopping rule for smooth local preprocessing functions $\phi$ beyond the explicit cases.

The assumptions delimit the contribution. We use one synchronized complex channel use, unit effective channel gains, the same encoder at every transmitter, and at least four transmitters. Unequal uncompensated gains, user-specific encoders, average-power constraints, and unrestricted $K=3$ classification require different arguments. The desired function $f:I^K\to\R$ is fixed before design; the result does not cover arbitrary symmetric functions. The extension in Section~\ref{sec:extension} requires common-completion identifiability for encoder classification, Lipschitz continuity for a positive design optimum, and separate endpoint conditions for the explicit quadratic formula. Universal sum representations have established dimension limitations~\cite{wagstaff19,wagstaff22}.

For a fixed encoder, the score determines the smallest uniform disturbance-amplification constant over all exact decoders. The constrained frontier minimizes the leading MSE coefficient of the specified strip decoder; finite-noise Bayesian optimization over arbitrary estimators remains a separate question. Section~\ref{sec:implementation} bounds residual mismatch without solving optimal design for unequal channels. Appendix~\ref{app:boundaries} gives precise boundaries for criterion transfer, poor conditioning, and finite alphabets. Appendix~\ref{app:approx-rigidity} proves qualitative stability under common regularity bounds, while an explicit error-to-geometry rate remains open. These restrictions identify concrete directions for extending the present classification and design theorem.
